% New English translation from checked cote 25, site 86 (PDF page 87).
% AI-translated by GPT-6 Astra (OpenAI, ChatGPT Pro mode).
\scanpage{25/86}
\mspage{\hand{\unsure{80}}}

p$_{\var{\deleted{\unsure{C}}}{}{deletion}{ega25_25_086-v0001}}$ above, \var{\deleted{which it will be convenient to denote in the same way as well}}{}{deletion}{ega25_25_086-v0002}
\var{\deleted{p$_C$, to indicate its dependence on C.}}{}{deletion}{ega25_25_086-v0003} In this way, the properties of
the restricted conical projection reduce \var{\deleted{to a large extent}}{}{deletion}{ega25_25_086-v0004} to
those of the extended conical projection, which has been studied systematically
elsewhere, or is supposed to have been (cf n\textsuperscript{o} 10 and n\textsuperscript{o} 12). The
main question\var{\xout{s}}{}{typewriter strike-out}{ega25_25_086-v0005} that arises, when S=Spec(k), is what are
the properties of the conical p\unsure{ro|oo}jection of X when we take C generic
in Grass$^{m+1}$(\unsure{E|P}) (which requires a base change
k$\hand{\rightarrow}$k(\hand{$\unsure{\eta|y}$}), i.e. C is \var{\inserted[typed]{then}}{then}{insertion}{ega25_25_086-v0007} in fact a linear subvariety of X$_{k(\hand{\unsure{\eta|y}})}$); standard
\unsure{ar|aa}guments, already repeated ad nauseam, will allow us to deduce
the analogous properties for conical projections corresponding to \var{\xout{of}}{}{typewriter strike-out}{ega25_25_086-v0008}
points of Grass$^{m+1}$(P) belonging to a nonempty open subset of that
Grassmannian, and finally, when k is infinite, we conclude that there exists
one (in fact infinitely many) \var{\inserted[typed]{of}}{of}{insertion}{ega25_25_086-v0009} \var{\xout{of}}{}{typewriter strike-out}{ega25_25_086-v0010} C \uline{defined} \uline{over} k, i.e. a linear
subvariety of P itself (without changing the base field), giving rise
to a conical projection with the properties in question. We should
\var{\margin{|}}{\placedmargin{|}}{marginal note placed}{ega25_25_086-v0012} combine general explanations of this sort with those of the same
kind given in n\textsuperscript{os} 4 and 7, which we have moreover already
used more or less implicitly, for example in n\textsuperscript{o}13 \var{\xout{\unsure{(proof of}}}{}{typewriter strike-out}{ega25_25_086-v0013}
\var{\xout{13.3}}{}{typewriter strike-out}{ega25_25_086-v0014} (cf 13.4 c). In this connection, we should also examine the
properties relative to a sheaf F on X, taking its inverse image
$F^{(m)}_{\hand{Q}(C)}$ on $\underline{X}^{(m)}_{\hand{Q}(C)}$. Moreover, in the precise situation
considered here, we need simpler notation; I propose $\overset{\hand{\sim}}{X}(C)$ and $\overset{\hand{\sim}}{F}(\unsure{C|O})$,
and simply $\overset{\hand{\sim}}{X}$ and $\overset{\hand{\sim}}{F}$ when \unsure{no} confusion results (note that the
F \var{\deleted{above}}{}{deletion}{ega25_25_086-v0015} \var{\inserted{\unsure{here}}}{\unsure{here}}{insertion}{ega25_25_086-v0016} is not the one at the beginning of the n\textsuperscript{o}).

Broadly speaking, and when, say, we assume f an immersion, the
properties of the generic conical projection are rather different,
\tnote{25-086-1}{The source describes C as a subvariety of the base-changed X; this letter is retained. The French adjective agreement corrections in principales and définies have no English effect. The source's doubtful apostrophe in n'en remains marked on the English negation no.}

\mspage{\hand{86}}
